% !PRINTED 641
\weberpage{641}

Die Quotienten ganzer Funktionen in $Z$ sind die Funktionale in $\overline{Z}$. Jedes Funktional $\overline{A}$ in $\overline{Z}$ kann in die Form gesetzt werden:
\weq{1}{\overline{A} = \frac{\Phi_1}{\Phi_2} = a\frac{E_1}{E_2} = aE,}
worin $\Phi_1$, $\Phi_2$, $E_1$, $E_2$ ganze Funktionen in $Z$ sind, darunter die Einheiten $E_1$, $E_2$, und $E$ ist eine Einheit in $\overline{Z}$, die sich als gebrochene Funktion darstellt. $a$ ist der Quotient der Teiler von $\Phi_1$ und $\Phi_2$, also eine ganze oder gebrochene rationale Funktion von $z$.

Die Funktionen $a$ und $E$ in der Formel (1) sind durch $\overline{A}$ völlig bestimmt, abgesehen von \gesperrt{konstanten Faktoren}, die hier die Rolle der \gesperrt{numerischen Einheiten} spielen. Wir wollen $a$ die \glqq \gesperrt{Absolute}\grqq{} von $\overline{A}$ nennen. Es gilt dann der Satz:

1. \gesperrt{Die Absolute eines Produktes ist gleich dem Produkt der Absoluten der Faktoren,}
und wir definieren:

2. \gesperrt{Ein rationales Funktional heißt ganz, wenn seine Absolute eine ganze Funktion von $z$ ist.}

Aus 1. folgt, daß das Produkt zweier ganzer Funktionale in $\overline{Z}$ wieder ein ganzes Funktional ist. Dasselbe folgt auch für die Summe und Differenz zweier ganzer Funktionale $\overline{A}_1$, $\overline{A}_2$ in $\overline{Z}$. Setzen wir nämlich:
\[
\overline{A}_1 = a_1\frac{E_1}{E}, \quad \overline{A}_2 = \frac{a_2E_2}{E},
\]
\[
\overline{A}_1 \pm \overline{A}_2 = \frac{a_1E_1 \pm a_2E_2}{E},
\]
worin $a_1$, $a_2$ ganze rationale Funktionen von $z$ sind; $E_1$, $E_2$, $E$ ganze Einheiten, so ist die Absolute von $\overline{A}_1 \pm \overline{A}_2$ der Teiler der ganzen Funktion $a_1E_1 \pm a_2E_2$, also auch eine ganze Funktion von $z$. Also gilt der Satz:

3. \gesperrt{Summe, Differenz und Produkt zweier ganzer Funktionale in $\overline{Z}$ sind wieder ganze rationale Funktionale.}

Alle Einheiten und deren Reziproken sind nach der Definition als \gesperrt{ganz} zu bezeichnen. Ist die Absolute eines ganzen rationalen Funktionals $\overline{A}$ linear ($= z - c$), so heißt $\overline{A}$ ein \gesperrt{rationales Primfunktional} oder ein \gesperrt{Primfunktional} in $\overline{Z}$.
